% GENERATED FILE. Edit canonical lessons, then run npm run generate:books.
% canonical-source-hash: fnv1a64:a3b057f30e9c09f5
% canonical-lessons: SA-C275-karisyati, SA-C275-pathisyati, SA-R275-what-everyone-will-do

\chapter{What Everyone Will Do}
\label{ch:sa-what-everyone-will-do}
\hlchaptermodality{275}

\begin{chapteropening}
\textbf{By the end of this chapter:} \emph{I can say what anyone will do, in every person and number (kariṣyāmi, kariṣyati, kariṣyanti), with each verb's own future stem (paṭhiṣyati, drakṣyati, dāsyati).}

Complete SA-R275-what-everyone-will-do: say what you, your friends and everyone else will do.
\end{chapteropening}

\section[\texorpdfstring{\sk{करिष्यति} (kariṣyati)}{करिष्यति (kariṣyati)}]{\sk{करिष्यति} (kariṣyati) — the future, person by person}

\label{lesson:SA-C275-karisyati}



\begin{quote}
Before the new one: say \textquotedblleft{}I will do\textquotedblright{}, \textbf{\sk{करिष्यामि}} (\emph{kariṣyāmi}), and \textquotedblleft{}I will go\textquotedblright{}, \textbf{\sk{गमिष्यामि}} (\emph{gamiṣyāmi}). Then say \textquotedblleft{}tomorrow\textquotedblright{}. (\emph{śvaḥ}.) This lesson gives that future to every person.
\end{quote}

\subsection*{You'll want to know: I, you, he and she}
\begin{quote}
\textbf{\sk{अहं} \sk{करिष्यामि।}}
\end{quote}

\begin{quote}
\emph{ahaṃ kariṣyāmi.} — \textquotedblleft{}I will do it.\textquotedblright{}
\end{quote}

\begin{quote}
\textbf{\sk{त्वं} \sk{करिष्यसि।}}
\end{quote}

\begin{quote}
\emph{tvaṃ kariṣyasi.} — \textquotedblleft{}You will do it.\textquotedblright{}
\end{quote}

\begin{quote}
\textbf{\sk{सः} \sk{करिष्यति।}}
\end{quote}

\begin{quote}
\emph{saḥ kariṣyati.} — \textquotedblleft{}He will do it.\textquotedblright{}
\end{quote}

\textbf{\sk{करिष्य}-} (\emph{kariṣya-}) is the future stem of \textquotedblleft{}do\textquotedblright{}, and the present's endings follow it: \emph{kariṣyāmi} like \emph{gacchāmi}, \emph{kariṣyati} like \emph{gacchati}.

\begin{quote}
\textbf{\sk{सः} \sk{श्वः} \sk{कार्यं} \sk{करिष्यति।}}
\end{quote}

\begin{quote}
\emph{saḥ śvaḥ kāryaṃ kariṣyati.} — \textquotedblleft{}He will do the work tomorrow.\textquotedblright{}
\end{quote}

\subsection*{You'll want to know: we, you all, they, and the two}
\textbf{\sk{वयम्}} (\emph{vayam}) is \textquotedblleft{}we\textquotedblright{}, \textbf{\sk{यूयम्}} (\emph{yūyam}) \textquotedblleft{}you all\textquotedblright{}, and \textbf{\sk{ते}} (\emph{te}) \textquotedblleft{}they\textquotedblright{}.

\begin{quote}
\textbf{\sk{वयं} \sk{करिष्यामः।}}
\end{quote}

\begin{quote}
\emph{vayaṃ kariṣyāmaḥ.} — \textquotedblleft{}We will do it.\textquotedblright{}
\end{quote}

\begin{quote}
\textbf{\sk{ते} \sk{करिष्यन्ति।}}
\end{quote}

\begin{quote}
\emph{te kariṣyanti.} — \textquotedblleft{}They will do it.\textquotedblright{}
\end{quote}

\noindent
\begin{tabularx}{\linewidth}{@{}>{\raggedright\arraybackslash}X>{\raggedright\arraybackslash}X@{}}
\toprule
\textbf{who} & \textbf{the future of \textquotedblleft{}do\textquotedblright{}} \\
\midrule
\textbf{\sk{अहम्}} (\emph{aham}), I & \textbf{\sk{करिष्यामि}} (\emph{kariṣyāmi}) \\
\textbf{\sk{त्वम्}} (\emph{tvam}), you & \textbf{\sk{करिष्यसि}} (\emph{kariṣyasi}) \\
\textbf{\sk{सः}}, \textbf{\sk{सा}} (\emph{saḥ}, \emph{sā}), he, she & \textbf{\sk{करिष्यति}} (\emph{kariṣyati}) \\
\textbf{\sk{वयम्}} (\emph{vayam}), we & \textbf{\sk{करिष्यामः}} (\emph{kariṣyāmaḥ}) \\
\textbf{\sk{यूयम्}} (\emph{yūyam}), you all & \textbf{\sk{करिष्यथ}} (\emph{kariṣyatha}) \\
\textbf{\sk{ते}} (\emph{te}), they & \textbf{\sk{करिष्यन्ति}} (\emph{kariṣyanti}) \\
\bottomrule
\end{tabularx}

Exactly two people take the dual, as \textbf{\sk{वदावः}} (\emph{vadāvaḥ}), \textquotedblleft{}we two speak\textquotedblright{}, did:

\begin{quote}
\textbf{\sk{अरुणः} \sk{मीरा} \sk{च} \sk{करिष्यतः।}}
\end{quote}

\begin{quote}
\emph{aruṇaḥ mīrā ca kariṣyataḥ.} — \textquotedblleft{}Arun and Mira will do it.\textquotedblright{}
\end{quote}

A question puts \textbf{\sk{किम्}} (\emph{kim}) first; \textquotedblleft{}not\textquotedblright{} puts \textbf{\sk{न}} (\emph{na}) before the verb:

\begin{quote}
\textbf{\sk{किं} \sk{त्वं} \sk{श्वः} \sk{गमिष्यसि}?}
\end{quote}

\begin{quote}
\emph{kiṃ tvaṃ śvaḥ gamiṣyasi?} — \textquotedblleft{}Will you go tomorrow?\textquotedblright{}
\end{quote}

\begin{quote}
\textbf{\sk{न}, \sk{अहं} \sk{न} \sk{गमिष्यामि।}}
\end{quote}

\begin{quote}
\emph{na, ahaṃ na gamiṣyāmi.} — \textquotedblleft{}No, I will not go.\textquotedblright{}
\end{quote}

\subsection*{Your turn}
Say these aloud:

\begin{itemize}
\raggedright
  \item \textquotedblleft{}I will go home\textquotedblright{} — \emph{ahaṃ gṛhaṃ gamiṣyāmi}
  \item \textquotedblleft{}We will go home\textquotedblright{} — \emph{vayaṃ gṛhaṃ gamiṣyāmaḥ}
  \item \textquotedblleft{}They will go home\textquotedblright{} — \emph{te gṛhaṃ gamiṣyanti}
\end{itemize}

\begin{tcolorbox}[breakable,colback=teal!4,colframe=teal!35!black,title={Before you move on}]
\textquotedblleft{}You will do it\textquotedblright{}? (\textbf{\sk{त्वं} \sk{करिष्यसि}}.) \textquotedblleft{}They will do it\textquotedblright{}? (\textbf{\sk{ते} \sk{करिष्यन्ति}}.) Arun and Mira, the two of them? (\textbf{\sk{करिष्यतः}}.) Which word opens the question? (\textbf{\sk{किम्}}.)
\end{tcolorbox}
\section[\texorpdfstring{\sk{पठिष्यति} · \sk{द्रक्ष्यति} (paṭhiṣyati …)}{पठिष्यति · द्रक्ष्यति (paṭhiṣyati …)}]{\sk{पठिष्यति} · \sk{द्रक्ष्यति} (paṭhiṣyati · drakṣyati) — each verb's own future}

\label{lesson:SA-C275-pathisyati}



\begin{quote}
Say \textquotedblleft{}she will do it\textquotedblright{} and \textquotedblleft{}we will do it\textquotedblright{}. (\emph{sā kariṣyati}, \emph{vayaṃ kariṣyāmaḥ}.) Then say \textquotedblleft{}I will read\textquotedblright{}. (\emph{paṭhiṣyāmi}.) The endings stay the same; the stem belongs to each verb.
\end{quote}

\subsection*{You'll want to know: -\sk{इष्य}-}
\begin{quote}
\textbf{\sk{सः} \sk{पठिष्यति।}}
\end{quote}

\begin{quote}
\emph{saḥ paṭhiṣyati.} — \textquotedblleft{}He will read.\textquotedblright{}
\end{quote}

Most verbs you know make the future with \textbf{-\sk{इष्य}-} (\emph{-iṣya-}) after the root: \textbf{\sk{पठति}} gives \textbf{\sk{पठिष्यति}}, and \textbf{\sk{खादति}} (\emph{khādati}), \textquotedblleft{}eats\textquotedblright{}, gives \textbf{\sk{खादिष्यति}} (\emph{khādiṣyati}).

\begin{quote}
\textbf{\sk{मीरा} \sk{पुस्तकं} \sk{पठिष्यति।}}
\end{quote}

\begin{quote}
\emph{mīrā pustakaṃ paṭhiṣyati.} — \textquotedblleft{}Mira will read the book.\textquotedblright{}
\end{quote}

\begin{quote}
\textbf{\sk{बालः} \sk{फलं} \sk{खादिष्यति।}}
\end{quote}

\begin{quote}
\emph{bālaḥ phalaṃ khādiṣyati.} — \textquotedblleft{}The boy will eat the fruit.\textquotedblright{}
\end{quote}

\textbf{\sk{वदति}} (\emph{vadati}), \textquotedblleft{}speaks\textquotedblright{}, gives \textbf{\sk{वदिष्यति}} (\emph{vadiṣyati}). \textbf{\sk{भवति}} gives \textbf{\sk{भविष्यति}} (\emph{bhaviṣyati}), the \textquotedblleft{}will be\textquotedblright{} you know. \textbf{\sk{लिखति}} (\emph{likhati}) changes its vowel: \textbf{\sk{लेखिष्यति}} (\emph{lekhiṣyati}), \textquotedblleft{}will write\textquotedblright{}.

\begin{quote}
\textbf{\sk{सा} \sk{श्वः} \sk{लेखिष्यति।}}
\end{quote}

\begin{quote}
\emph{sā śvaḥ lekhiṣyati.} — \textquotedblleft{}She will write tomorrow.\textquotedblright{}
\end{quote}

\subsection*{You'll want to know: -\sk{स्य}-}
A few verbs take \textbf{-\sk{स्य}-} (\emph{-sya-}) instead, on the root itself. \textbf{\sk{पश्यति}} (\emph{paśyati}), \textquotedblleft{}sees\textquotedblright{}, has the root √dṛś, and its future is \textbf{\sk{द्रक्ष्यति}} (\emph{drakṣyati}):

\begin{quote}
\textbf{\sk{सः} \sk{नगरं} \sk{द्रक्ष्यति।}}
\end{quote}

\begin{quote}
\emph{saḥ nagaraṃ drakṣyati.} — \textquotedblleft{}He will see the city.\textquotedblright{}
\end{quote}

\textbf{\sk{ददाति}} (\emph{dadāti}), \textquotedblleft{}gives\textquotedblright{}, has \textbf{\sk{दास्यति}} (\emph{dāsyati}); \textbf{\sk{तिष्ठति}} (\emph{tiṣṭhati}), \textquotedblleft{}stands, stays\textquotedblright{}, has \textbf{\sk{स्थास्यति}} (\emph{sthāsyati}):

\begin{quote}
\textbf{\sk{सा} \sk{पानीयं} \sk{दास्यति।}}
\end{quote}

\begin{quote}
\emph{sā pānīyaṃ dāsyati.} — \textquotedblleft{}She will give water.\textquotedblright{}
\end{quote}

\begin{quote}
\textbf{\sk{अहम्} \sk{अत्र} \sk{स्थास्यामि।}}
\end{quote}

\begin{quote}
\emph{aham atra sthāsyāmi.} — \textquotedblleft{}I will stay here.\textquotedblright{}
\end{quote}

\noindent
\begin{tabularx}{\linewidth}{@{}>{\raggedright\arraybackslash}X>{\raggedright\arraybackslash}X@{}}
\toprule
\textbf{the present} & \textbf{the future} \\
\midrule
\textbf{\sk{पठति}}, \textbf{\sk{खादति}}, \textbf{\sk{वदति}} & \textbf{\sk{पठिष्यति}}, \textbf{\sk{खादिष्यति}}, \textbf{\sk{वदिष्यति}} \\
\textbf{\sk{लिखति}} & \textbf{\sk{लेखिष्यति}} \\
\textbf{\sk{गच्छति}} & \textbf{\sk{गमिष्यति}} \\
\textbf{\sk{पश्यति}} & \textbf{\sk{द्रक्ष्यति}} \\
\textbf{\sk{ददाति}}, \textbf{\sk{तिष्ठति}} & \textbf{\sk{दास्यति}}, \textbf{\sk{स्थास्यति}} \\
\bottomrule
\end{tabularx}

So a future is learned with its verb, as a dictionary gives it: Monier-Williams lists \textbf{\sk{द्रक्ष्यति}} under √dṛś, and Lanman's reader gives each of these with its root.

\subsection*{Your turn}
Say these aloud:

\begin{itemize}
\raggedright
  \item \textquotedblleft{}He will read the book\textquotedblright{} — \emph{saḥ pustakaṃ paṭhiṣyati}
  \item \textquotedblleft{}He will see the book\textquotedblright{} — \emph{saḥ pustakaṃ drakṣyati}
  \item \textquotedblleft{}He will give the book\textquotedblright{} — \emph{saḥ pustakaṃ dāsyati}
\end{itemize}

\begin{tcolorbox}[breakable,colback=teal!4,colframe=teal!35!black,title={Before you move on}]
\textquotedblleft{}Will eat\textquotedblright{}? (\textbf{\sk{खादिष्यति}}.) \textquotedblleft{}Will see\textquotedblright{}? (\textbf{\sk{द्रक्ष्यति}}: the stem comes with the verb.) \textquotedblleft{}Will give\textquotedblright{}? (\textbf{\sk{दास्यति}}.) And \textquotedblleft{}I will stay\textquotedblright{}? (\textbf{\sk{स्थास्यामि}}.)
\end{tcolorbox}
\section[(review)]{Review — what everyone will do}

\label{lesson:SA-R275-what-everyone-will-do}



\begin{quote}
Say \textquotedblleft{}I will do it\textquotedblright{} and \textquotedblleft{}he will read\textquotedblright{}. (\emph{ahaṃ kariṣyāmi}, \emph{saḥ paṭhiṣyati}.)
\end{quote}

\subsection*{Your turn}
Say these aloud:

\begin{itemize}
\raggedright
  \item Tomorrow I will read a book.
  \item Will you go home? No, I will not go.
  \item Arun and Mira will see the city.
  \item We will eat the fruit, and they will give water.
  \item She will stay here tomorrow.
\end{itemize}

\begin{tcolorbox}[breakable,colback=teal!4,colframe=teal!35!black,title={Before you move on}]
\textbf{\sk{करिष्यन्ति}}: how many people? (Three or more.) \textbf{\sk{करिष्यतः}}? (Two.) Now the other way round: \textquotedblleft{}will see\textquotedblright{}? (\textbf{\sk{द्रक्ष्यति}}.) \textquotedblleft{}Will stay\textquotedblright{}? (\textbf{\sk{स्थास्यति}}.) \textquotedblleft{}We will write\textquotedblright{}? (\textbf{\sk{वयं} \sk{लेखिष्यामः}}.)
\end{tcolorbox}
